% Experimental setup — draft. Do not mention the course assignment or its grading.
% Inference conditions are fixed by design; implementation in src/ktide/.
\section{Experimental Setup}
\label{sec:setup}

\subsection{Model and inference conditions}
\label{sec:setup-inference}

All models start from \texttt{Qwen2.5-1.5B-Instruct} \todo{cite Qwen2.5 technical report after verification}, and the architecture is unchanged.
Inference conditions are fixed for every model: FP16 on a CUDA GPU, SDPA attention, greedy decoding (no sampling, one beam), at most 128 new tokens, batch size 4 with left padding, and one fixed Korean system prompt that states the answer constraints.
The system prompt is never tuned per model.
Every evaluated model is a full, non-quantized checkpoint with the original tokenizer and chat template.
Before each evaluation, we compute the effective generation configuration to confirm greedy decoding (Appendix~\ref{app:genconfig}).

\subsection{Models compared}
\label{sec:setup-design}

\begin{table}[t]
\centering
\small
\begin{tabular}{lll}
\toprule
 & Translated terms & Preserved English terms \\
\midrule
Lecture-only & M1 & M2 \\
Enriched (CKU) & M3 & M4 \\
\bottomrule
\end{tabular}
\caption{The $2\times2$ design. M0 is the base model without fine-tuning.}
\label{tab:design}
\end{table}

The main study is a $2\times2$ factorial design (Table~\ref{tab:design}) plus the base model M0.
Comparing M1 with M2 and M3 with M4 isolates terminology preservation (RQ2).
Comparing M1 with M3 and M2 with M4 isolates knowledge enrichment (RQ1), and the four cells together estimate their interaction.
All four fine-tuned models share the teacher model, concept coverage, question-type distribution, number of training examples, LoRA configuration and number of optimizer steps \pending{exact values}.
\pending{Seeds per condition; any condition that could not be controlled.}
\todo{Optional if time allows: full fine-tuning vs.\ LoRA on the best cell; one additional small model from another family.}

\subsection{Evaluation sets}
\label{sec:setup-data}

\pending{Independent evaluation set built and frozen before any training data is generated: size; stratification by domain $\times$ question type $\times$ term surface form (Korean term, English term, acronym); development/held-out split; a seen-concept paraphrase subset and an unseen-concept subset to distinguish memorization from generalization; the rule that held-out items never enter teacher prompts.}

\subsection{Metrics}
\label{sec:setup-metrics}

\paragraph{Rubric-based LLM judge.}
Each answer is scored on accuracy, relevance and constraint compliance (1--5 each) from the question and answer alone.
The judge is \texttt{gpt-6.2-sol} with medium reasoning effort, run through the Codex CLI in an empty working directory with no project context.
Answers are judged in batches of 20 per call; averages are recomputed from the sub-scores, and the raw judge output is kept.
Answers are anonymized, so the judge never sees the model name or training condition.

\paragraph{Judge variance and human agreement.}
In a pilot, judging the same answers twice changed at least one sub-score for most items.
We therefore judge each run several times and report the mean and standard deviation \pending{number of repetitions}.
We also measure agreement with human ratings on a subset \pending{size and agreement statistic}.

\paragraph{Rule-based metrics.}
Computed in Python without an LLM:
word-limit compliance (at most 30 whitespace-separated words),
sentence-limit compliance (one or two sentences, with a documented splitting rule),
empty answers, answers that hit the token limit,
and unintended script mixing (e.g., Chinese characters).

\paragraph{Reference-based and terminology metrics.}
\pending{Reference-based semantic accuracy against verified CKU reference answers; required-concept coverage;
terminology preservation rate (TPR) and acronym expansion accuracy (AEA), which we define for this study and do not claim as standard metrics; hallucination rate.}

\paragraph{Uncertainty.}
\pending{Paired bootstrap confidence intervals over evaluation items for every main comparison; interaction test for the $2\times2$ design.}

\subsection{Training and compute}
\label{sec:setup-compute}

Each condition is trained with LoRA on all linear projections of the attention and MLP blocks (rank 16, $\alpha=32$, dropout 0.05), learning rate $10^{-4}$ with cosine decay and 3\% warm-up, three epochs, effective batch size 32, maximum length 256 tokens, and seed 42.
The loss is computed on the assistant answer and its end-of-turn token only; the system prompt and question are masked.
Five percent of the training examples are held out as a validation set for monitoring.
The adapter is merged into the base weights in FP32, the merged model is saved in FP16, and the original generation configuration is restored (Appendix~\ref{app:genconfig}).

All controlled runs are trained in FP16 mixed precision on a single GPU type (NVIDIA RTX 5000 Ada Generation) with FP32 adapter weights.
We log the gradient-scaler scale and count non-finite loss or gradient-norm values in every run.
For M1, training takes 27 optimizer steps (18.6\,s), and the run recorded no non-finite values and no reduction of the gradient scale.
An earlier M1 run, trained in BF16 on a different GPU before the precision was fixed, is kept as a pilot and excluded from the main comparison (Appendix~\ref{app:precision}).
\pending{Run times and GPU hours for M2--M4; judge cost.}
